\section{Contexto del Caso de Estudio}

\textbf{AbejaNet Revival} es una plataforma IoT para monitoreo en tiempo real de colmenas
de abejas (\textit{Apis mellifera}). El sistema integra sensores ESP32 instalados en las
colmenas para recopilar variables ambientales (temperatura, humedad, peso, decibelios y
lluvia), un backend Node.js/Express con base de datos PostgreSQL, una aplicaci\'on m\'ovil
React Native/Expo y un panel web de administraci\'on.

La tabla \texttt{lecturas\_ambientales} de la base de datos operacional constituye la
fuente principal de datos de esta evaluaci\'on. Dado que el servidor de producci\'on
(Render) no estaba activo durante el periodo de desarrollo, los datos se generaron de forma
sint\'etica con el script \texttt{00\_generar\_datos.py}, replicando la l\'ogica de
simulaci\'on del archivo original \texttt{backend/generate\_mock\_data.js} del proyecto.

\subsection{Arquitectura del Sistema}

\begin{itemize}\setlength{\itemsep}{2pt}
  \item \textbf{Sensores ESP32:} Env\'ian lecturas v\'ia POST \texttt{/api/sensor-data}
        cada 10 minutos con autenticaci\'on por clave API.
  \item \textbf{Backend Node.js/Express:} Recibe, valida y persiste los datos en
        PostgreSQL alojado en Render.
  \item \textbf{Aplicaci\'on m\'ovil:} Consume la API REST con autenticaci\'on JWT,
        mostrando gr\'aficas de tendencias por colmena.
  \item \textbf{Panel web:} Gesti\'on de usuarios, apiarios, colmenas y alertas.
\end{itemize}

\subsection{Colmenas Monitoreadas}

\begin{table}[h!]
  \centering
  \rowcolors{2}{grisCelda}{white}
  \begin{tabular}{llll}
    \toprule
    \textbf{Apiario} & \textbf{Colmena} & \textbf{MAC} & \textbf{Tipo Sensor} \\
    \midrule
    Apiario Principal   & Colmena Alfa Ppal  & AA:BB:CC:11:22:33 & Temperatura/Humedad \\
    Apiario Principal   & Colmena Gamma Ppal & AA:BB:CC:11:22:44 & Multisensor \\
    Apiario Laboratorio & Colmena Beta Lab   & A2:04:2A:B9:C1:D9 & Multisensor \\
    \bottomrule
  \end{tabular}
  \caption{Inventario de colmenas y sensores del sistema AbejaNet}
\end{table}
